\section{F-SM-29: the underlying connective localization over every scalar ring}

\subsection{Precise statement and distinction from the ordinary right adjoint}
Let \(k\) be an arbitrary commutative ring and \(R\) a small
\(k\)-linear category.  Set
\(\mathcal M=\operatorname{Fun}_k(R^{\mathrm{op}},\operatorname{Mod}k)\).
Assume the following data from BR-13403:
\begin{enumerate}
\item an exact scalar functor \(\kappa(U)=T\otimes_kU\), preserving
coproducts, with left adjoint \(\varphi\);
\item an exact additive left adjoint
\(\Lambda:\mathcal M\to\operatorname{Ch}_{\ge0}(k)\), with a right
adjoint \(G\), such that \(H_0\Lambda=\varphi\) and
\(H_j\Lambda(P)=0\) for \(j>0\) and projective \(P\);
\item the canonical map
\(\Lambda\kappa(k)\to k[0]\), obtained from the counit of
\(\varphi\dashv\kappa\), is a quasi-isomorphism;
\item for every \(X\) there is the interval monomorphism
\(X\to J\otimes X\), and
\(L\varphi(J\otimes X)=0\), including degree zero.
\end{enumerate}
Put \(W=\Lambda^{-1}(\mathrm{qis})\).  Then the functor induced by
\(\Lambda\), equivalently by \(L\varphi\), is an equivalence
\[
 W^{-1}\mathcal M\xrightarrow{\ \sim\ }D_{\ge0}(k).
 \tag{SM29.1}
\]
There is no restriction on the cardinality or homological dimensions
of \(k\), no finite-generation restriction on the scalar morphism
modules, and no \(\mathrm{FP}_\infty\) hypothesis on \(D\).
In particular BR-13405 supplies these exact-adjunction data under
the corollary's stated locally projective ringoid hypotheses.

The source identification here is precise.  On manuscript footer 47,
condition (a) supplies the left adjoint \(\varphi\) of \(\kappa\)
and identifies its right adjoint as
\(X\mapsto\operatorname{Hom}_{\mathcal M}(\mathbf1,X)\).
Thus exactness and preservation of coproducts of \(\kappa\) follow
from its two adjoints; footer 46 also explicitly writes
``\(\kappa\) exact.''  Condition (b) displays the derived adjunction,
and (c) displays \(L\varphi(\kappa(k))\to k\) as an isomorphism.
We use this arrow as the displayed adjunction's counit, not just an
unspecified isomorphism between these objects.  Footer 48 adds the
finitely presented projective generators, and its right margin adds
scalar projectivity of the morphism modules.  These are the
module-category corollary's hypotheses.  We do not assert that the
broader abstract conditions (a)--(d), without this presentability
and exact-adjunction construction, have been settled by this proof.

The inverse constructed below uses a homotopy colimit.
The proof does not assert that the ordinary right adjoint \(G\)
preserves every quasi-isomorphism, nor that
\(\Lambda G(C)\to C\) is a quasi-isomorphism for every ordinary
complex \(C\).  That stronger assertion remains precisely the
acyclic-projective test of BR-13403.  The right \emph{derived}
adjoint is a quasi-inverse by (SM29.1).

\subsection{Exact background inputs}
We use Smith's model-structure recognition theorem as proved in
Beke (2000), Theorem 1.7, author-PDF p.~8.
Its hypotheses are checked below.  The accessibility facts used
are Beke's Proposition 1.15, p.~12, and Proposition 1.18, p.~13.
For a cofibrantly generated model category, the projective model
structure on diagrams over a small category is the standard
diagram-model theorem: see Dugger, current author version,
Theorem 19.4, p.~91, and the discussion of its generating
cofibrations in Beke, author-PDF p.~20.
We use Dold--Kan for abelian categories in the precise form of
the Stacks Project, Theorem 14.24.3, Tag \texttt{019G}.

The homotopy-colimit totalization formula needed in the argument
is proved below, including its naturality and augmentation
compatibility.  It also agrees with Dugger's Proposition 19.9,
p.~93 of the current 109-page author PDF.
None of these known background theorems is attributed to Alp\"oge.

\subsection{Complete proof}

\paragraph{1. All scalar modules satisfy canonical normalization.}
Write
\[
 e_U:\Lambda\kappa(U)\longrightarrow U[0]
 \tag{SM29.2}
\]
for the natural composite through
\(H_0\Lambda\kappa(U)[0]=\varphi\kappa(U)[0]\) and the adjunction
counit.  The hypothesis says \(e_k\) is a quasi-isomorphism.
Both \(\Lambda\) and \(\kappa\) preserve coproducts.
Consequently \(e_U\) is a quasi-isomorphism for free \(U\), and
then for projective \(U\) by retracts.

For arbitrary \(U\), take a projective scalar resolution
\(Q_\bullet\to U\) in nonnegative homological degrees.
Apply the exact functor \(\Lambda\kappa\) to obtain the
first-quadrant double complex
\[
 E_{a,b}=(\Lambda\kappa Q_a)_b.
\]
The augmentation
\(\operatorname{Tot}E\to\Lambda\kappa U\) is a quasi-isomorphism:
in each internal degree \(b\), the augmented horizontal sequence
is exact.  The maps \(e_{Q_a}\) give a quasi-isomorphism
\(\operatorname{Tot}E\to Q_\bullet\), by taking vertical homology.
The augmentation \(Q_\bullet\to U[0]\) is a quasi-isomorphism.
These maps form a commutative square with \(e_U\), so \(e_U\)
is a quasi-isomorphism.  Each total degree has finitely many
summands; both filtration arguments therefore have finite
diagonals.  We have proved the \emph{canonical}, natural assertion
\[
 \Lambda\kappa(U)\simeq U[0]\quad\text{for every scalar }U.
 \tag{SM29.3}
\]

\paragraph{2. The model structure on \(\mathcal M\).}
The module category of a small ringoid is locally finitely
presentable.  Concretely, presentations use coproducts of the
representables \(P_x=R(-,x)\), and every presentation is the filtered
colimit of its finite sets of generators and relations.  Limits,
colimits and exactness are computed objectwise.
In particular \(\mathcal M\) is a Grothendieck abelian category.

There is an explicit set of generating monomorphisms:
\[
 \mathcal I=\{\,U\hookrightarrow P_x:
       x\in R,\ U\text{ a submodule of }P_x\,\}.
 \tag{SM29.4}
\]
It is a set, because \(R\) is small and each representable has a
set of subobjects.  We verify that every monomorphism is a
transfinite composition of pushouts of maps in \(\mathcal I\).
Identify a given monomorphism with \(A\subset B\), and well-order
all the elements of all modules \(B(x)\).  At a successor stage,
adjoin the next element \(b\in B(x)\) to the submodule \(A_\beta\)
already constructed.  It defines \(P_x\to B\).  Let
\(U\subset P_x\) be the inverse image of \(A_\beta\).
Then the pushout
\[
 A_\beta\amalg_U P_x
\]
maps isomorphically onto \(A_\beta+\operatorname{im}(P_x\to B)\):
the kernel of \(P_x\to B\) is contained in \(U\), so the pushout
imposes exactly the needed identifications.  At a limit stage
take the union.  The construction eventually gives \(B\).
Conversely pushouts, transfinite compositions and retracts of
monomorphisms remain monomorphisms in modules.  Thus
\(\operatorname{cof}(\mathcal I)\) is exactly the monomorphisms.

The weak-equivalence class \(W\) is accessible and accessibly
embedded in the arrow category.  Indeed, \(\Lambda\), being a
left adjoint, preserves colimits; homology of module complexes
commutes with filtered colimits.  Hence
\[
 H_*\Lambda:\mathcal M\longrightarrow
       \prod_{n\ge0}\operatorname{Mod}k
\]
is accessible.  The class of isomorphisms in the arrow category
of graded modules is accessible: evaluation on the source is an
equivalence from that class to graded modules, and its inclusion
preserves filtered colimits.  The accessible inverse-image
criterion, Beke Proposition 1.18, applies to \(H_*\Lambda\).
Proposition 1.15 then gives Smith's solution-set condition at
the set \(\mathcal I\).

We check the other recognition conditions individually.
Closure under retracts and two-out-of-three follow from the
same properties of quasi-isomorphisms.
A monomorphism lies in \(W\) if and only if its cokernel has
acyclic \(\Lambda\)-image, because \(\Lambda\) is exact.
A pushout of a monomorphism has the same cokernel, so pushouts
preserve monomorphisms in \(W\).
Transfinite compositions also preserve this class: \(\Lambda\)
preserves the relevant colimits, filtered colimits of modules
are exact, and all the induced maps on homology in the
transfinite sequence are isomorphisms.

Finally, a map with the right lifting property against
\(\mathcal I\) has it against every monomorphism.
Lifting against \(0\to Y\), where \(Y\) is its target, supplies
a section.  Its kernel \(I\) is injective: a map from a subobject
to \(I\) extends across the containing object by applying the
same lifting property with zero map to \(Y\).
Thus the original map is a split epimorphism with injective
kernel.  The interval monomorphism
\(I\to J\otimes I\) splits, by injectivity, so
\[
 L\varphi(I)=0
\]
as a retract of \(L\varphi(J\otimes I)=0\).
BR-13401 identifies \(L\varphi(I)\) with \(\Lambda I\).
The split epimorphism consequently lies in \(W\).

All four conditions of Beke Theorem 1.7 are now verified.
There is a cofibrantly generated model structure on
\(\mathcal M\) with weak equivalences \(W\) and all monomorphisms
as cofibrations.  Every object is cofibrant.  Its homotopy
category is the localization \(W^{-1}\mathcal M\).

The same recognition argument gives a model structure on
\(\operatorname{Ch}_{\ge0}(k)\) with quasi-isomorphisms as weak
equivalences and all monomorphisms as cofibrations.
This category is itself a module category: its finitely presented
projective generators are \(k[0]\) and the disks \(D^n(k)\),
\(n\ge1\), by the evaluation adjunctions in BR-13402.
For its last recognition condition, an injective chain object
embeds into \(D^1(k)\otimes I\), a contractible complex, via
the degree-zero endpoint; the embedding splits and therefore
the injective chain object is acyclic.  All the other
conditions are exactly the preceding module-category argument.
This construction is within nonnegative complexes; no
unbounded-model restriction is being assumed.
The target homotopy category is \(D_{\ge0}(k)\).

Exactness of \(\Lambda\) and the definition of \(W\) show that
\(\Lambda\) preserves cofibrations and trivial cofibrations.
Together with its given right adjoint \(G\), it is therefore
a Quillen adjunction.  Its left derived functor is represented
by \(\Lambda\) itself, since every source object is cofibrant.

\paragraph{3. Homotopy colimits and the normalized total complex.}
We record the precise realization fact used below.
For a simplicial object
\(Y:\Delta^{\mathrm{op}}\to\operatorname{Ch}_{\ge0}(k)\),
there is a natural isomorphism in \(D_{\ge0}(k)\)
\[
 \operatorname*{hocolim}_{\Delta^{\mathrm{op}}}Y
       \simeq\operatorname{Tot}N_{\mathrm{simp}}Y.
 \tag{SM29.5}
\]
The totalization is the direct-sum totalization, with the usual
sign on one differential.  Its diagonals are finite.

Here is a proof in the model structures just specified.
Use the projective model structure on simplicial diagrams and
a functorial projectively cofibrant replacement \(QY\to Y\).
The colimit/constant-diagram adjunction is Quillen, since the
constant functor preserves the objectwise fibrations and trivial
fibrations of the projective diagram structure.
By definition the homotopy colimit is \(\operatorname{colim}QY\).
There is a natural augmentation
\[
 \theta_Z:\operatorname{Tot}N_{\mathrm{simp}}Z
              \longrightarrow\operatorname{colim}Z
\]
for every such diagram \(Z\), induced by its degree-zero
augmentation and zero on positive simplicial degrees.

Let \(F_{[m]}\) be left adjoint to evaluation at \([m]\) on
simplicial diagrams.  It is the copower by the representable
simplicial set \(\Delta[m]\).  For every chain complex \(C\),
\[
 \operatorname{Tot}N_{\mathrm{simp}}F_{[m]}(C)
       =N(k[\Delta[m]])\otimes_k C,\qquad
 \operatorname{colim}F_{[m]}(C)=C.
\]
The normalized chains of \(\Delta[m]\) contract onto \(k[0]\).
Explicitly, on a nondegenerate ordered simplex prepend the
vertex \(0\); if this creates a repeated vertex, its normalized
value is zero.  The boundary identity is
\(dh+hd=1-\iota\epsilon\).
Tensoring this contraction with \(C\), with the standard total
sign, proves that \(\theta_{F_{[m]}C}\) is a quasi-isomorphism
for arbitrary \(C\).

Projective diagram cofibrations are generated by the maps
\(F_{[m]}(i)\), with \(i\) in a set of generating target
cofibrations.  Both normalized totalization and colimit
preserve colimits; normalization does so by Dold--Kan,
and totalization does so by direct sums.
For a pushout adjoining a cell \(F_{[m]}(A)\to F_{[m]}(B)\),
both resulting maps from the preceding stage are monomorphisms.
Indeed, before the pushout these maps are the tensor of \(A\to B\)
with the degreewise free complex \(N(k[\Delta[m]])\), and
\(A\to B\) itself, respectively.
Their cokernels are respectively
\(N(k[\Delta[m]])\otimes(B/A)\) and \(B/A\).
The comparison on those cokernels is a quasi-isomorphism by
the preceding contraction.  The long exact homology sequences
therefore propagate the assertion about \(\theta\) across the
cell pushout.  Coproducts of cells, transfinite compositions
and retracts preserve it, since direct sums and filtered
colimits of modules are exact.  Every projectively cofibrant
diagram is a retract of such a cell diagram.  Hence
\(\theta_{QY}\) is a quasi-isomorphism.

The objectwise quasi-isomorphism \(QY\to Y\) also induces a
quasi-isomorphism on normalized total complexes.
Indeed normalization is exact in the simplicial direction,
so internal homology commutes with it.  The vertical-homology
comparison and the first-quadrant filtration prove the claim.
We have the natural zigzag
\[
 \operatorname{colim}QY
 \xleftarrow{\ \sim\ }\operatorname{Tot}N_{\mathrm{simp}}QY
 \xrightarrow{\ \sim\ }\operatorname{Tot}N_{\mathrm{simp}}Y,
\]
which proves (SM29.5).  All maps commute with an augmentation
of \(Y\) to a constant diagram; this gives the augmentation
compatibility needed later.

For any small diagram shape \(I\), the left Quillen functor
\(\Lambda\) preserves homotopy colimits in the derived sense:
\[
 \Lambda\operatorname*{hocolim}_I Z
       \simeq\operatorname*{hocolim}_I\Lambda Z.
 \tag{SM29.6}
\]
To verify this precisely, use projective model structures on
the diagram categories.  The objectwise adjunction is
Quillen, since its right adjoint preserves the objectwise
fibrations and trivial fibrations.  Thus \(\Lambda\) sends a
projectively cofibrant replacement to a projectively
cofibrant diagram, and it preserves the objectwise weak
equivalences involved.  It commutes with ordinary colimit
because it is a left adjoint.  Computing both derived
colimits with these cofibrant diagrams gives (SM29.6),
naturally in \(Z\).

\paragraph{4. Construction of a homotopical inverse.}
Let \(\Gamma\) be the Dold--Kan inverse from nonnegative
scalar chain complexes to simplicial scalar modules.
Thus \(N\Gamma C=C\), and every \((\Gamma C)_n\) is a
finite direct sum of the modules \(C_j\), \(0\le j\le n\).
Choose functorial projective-diagram cofibrant replacements
in \(\mathcal M^{\Delta^{\mathrm{op}}}\), and let \(h\) be
the resulting ordinary functor representing homotopy colimit.
Define an ordinary functor on chain complexes by
\[
 \mathcal L(C)=h\bigl(\kappa\Gamma C\bigr)\in\mathcal M.
 \tag{SM29.7}
\]
Equations (SM29.3), (SM29.5), and (SM29.6) give a natural
isomorphism in the derived category
\[
 \begin{split}
 \Lambda\mathcal L(C)
 &\simeq\operatorname*{hocolim}_{\Delta^{\mathrm{op}}}
               \Lambda\kappa\Gamma C\\
 &\simeq\operatorname*{hocolim}_{\Delta^{\mathrm{op}}}
               (\Gamma C)[0]\\
 &\simeq N\Gamma C=C.
 \end{split}
 \tag{SM29.8}
\]
The middle comparison is induced by the actual natural maps
\(e_U\), which are quasi-isomorphisms for every coefficient
module.  No projectivity of the modules in \(C\) is required.

If \(f:C\to C'\) is a quasi-isomorphism, naturality of
(SM29.8) shows that \(\Lambda\mathcal L(f)\) is an
isomorphism in the derived category, hence is a
quasi-isomorphism of complexes.
By the definition of \(W\), \(\mathcal L(f)\in W\).
Thus \(\mathcal L\) descends to a functor
\(D_{\ge0}(k)\to W^{-1}\mathcal M\), with
\(\Lambda\mathcal L\cong1\).
This argument proves homotopicality of \(\mathcal L\);
it does not incorrectly assume that \(\kappa\Gamma\)
takes quasi-isomorphisms to objectwise weak equivalences.

\paragraph{5. The other natural isomorphism.}
There is a functorial projective resolution \(P_\bullet X\to X\)
in the module category \(\mathcal M\).
One explicit choice starts with the coproduct of \(P_x\)
indexed by every element of every \(X(x)\), maps it onto
\(X\), and iterates the same construction on its kernel.
Maps of modules send the indexing elements to their images,
so the construction is functorial.

For a projective object \(P\), the adjunction unit
\[
 \eta_P:P\longrightarrow\kappa\varphi P
\]
belongs to \(W\).  The source \(\Lambda P\) has no positive
homology.  The target has no positive homology by (SM29.3).
On zeroth homology the map is
\(\varphi\eta_P:\varphi P\to\varphi\kappa\varphi P\);
the triangle identity gives
\(\epsilon_{\varphi P}\varphi\eta_P=1\).
The map \(\epsilon_{\varphi P}\) is an isomorphism by
(SM29.3), so \(\varphi\eta_P\) is an isomorphism too.

Apply Dold--Kan inside the abelian category \(\mathcal M\)
to \(P_\bullet X\), and put
\[
 A_X=h\Gamma(P_\bullet X).
\]
Every simplicial term \(\Gamma(P_\bullet X)_n\) is a finite
sum of projectives.  Their units consequently give an
objectwise \(W\)-equivalence of simplicial diagrams
\[
 \Gamma(P_\bullet X)
       \longrightarrow\kappa\Gamma(\varphi P_\bullet X).
\]
Here \(\varphi\) and \(\kappa\) commute with the finite sums
and differential matrices defining \(\Gamma\).
Taking homotopy colimits gives a natural \(W\)-equivalence
\[
 A_X\longrightarrow\mathcal L(\varphi P_\bullet X).
 \tag{SM29.9}
\]

The augmentation of the projective resolution gives a natural
map \(A_X\to X\).  Applying \(\Lambda\) and (SM29.5)--(SM29.6)
identifies its derived image with
\[
 \operatorname{Tot}_{a,b}(\Lambda P_aX)_b
               \longrightarrow\Lambda X.
 \tag{SM29.10}
\]
The identification is compatible with the augmentation by
Step 3.  Additivity of \(\Lambda\) makes it commute with the
finite-sum formula for \(\Gamma\), so normalizing that simplicial
direction gives \(\Lambda P_aX\) in degree \(a\).
For every \(b\), the augmented horizontal sequence
is exact because \(\Lambda\) is exact.
The first-quadrant filtration shows that (SM29.10) is a
quasi-isomorphism.  Therefore \(A_X\to X\) belongs to \(W\).

Finally define
\(T_X=\operatorname{Tot}_{a,b}(\Lambda P_aX)_b\).
The natural augmentations already used in BR-13401 give
\[
 \varphi P_\bullet X
       \xleftarrow{\ \sim\ }T_X
       \xrightarrow{\ \sim\ }\Lambda X.
 \tag{SM29.11}
\]
The left arrow is a quasi-isomorphism by the positive-degree
acyclicity of \(\Lambda\) on projectives, and the right arrow
is (SM29.10).  All comparisons are natural in \(X\).
Since \(\mathcal L\) is homotopical by Step 4, the zigzag
\[
 X\xleftarrow{\ \sim\ }A_X
   \xrightarrow{\ \sim\ }\mathcal L(\varphi P_\bullet X)
   \xleftarrow{\ \sim\ }\mathcal L(T_X)
   \xrightarrow{\ \sim\ }\mathcal L(\Lambda X)
 \tag{SM29.12}
\]
consists of maps in \(W\).
After localization it supplies
\(1\cong\mathcal L\Lambda\), naturally.
Together with (SM29.8), this proves (SM29.1).
Every bicomplex in the proof lies in the first quadrant;
no infinite-diagonal convergence claim has entered.
\(\square\)

\subsection{What this settles and what remains separate}
This establishes the original \emph{underlying} localization
conclusion for the finitely presented-projective-generator
corollary once the local scalar projectivity supplies the
exact adjunction BR-13405.  It works over every scalar ring.
The proof only needs the interval to annihilate injective
objects for the Smith recognition step; its supplied full
hypothesis guarantees this.

The distinction concerning \(G\) is substantive.
A right Quillen functor is applied to a fibrant replacement
when its right derived functor is formed.
The equivalence just proved therefore provides
\(\Lambda G(C^{\mathrm{fib}})\simeq C^{\mathrm{fib}}\)
at that derived level, with the usual cofibrant conventions,
but does not eliminate the test (134.6) for every ordinary
acyclic projective chain object \(P\).
F-SM-24 and F-SM-27 remain positive results for that stronger
ordinary-adjoint assertion.
The general handwritten theorem without the module-category
or comparable presentability assumptions is not claimed here.

\paragraph{References and version-specific locators.}
Beke, T. (2000). Sheafifiable homotopy model categories.
\emph{Mathematical Proceedings of the Cambridge Philosophical
Society, 129}(3), 447--475.
\url{https://doi.org/10.1017/S0305004100004722}.
Author PDF:
\url{https://faculty.uml.edu/tbeke/sheafi.pdf}.
Theorem 1.7, p.~8; Propositions 1.15 and 1.18, pp.~12--13;
diagram-model discussion p.~20.  These are author-PDF pages,
not asserted journal-page translations.

Dugger, D. (2008). \emph{A primer on homotopy colimits}
[Unpublished notes; current 109-page author version retrieved
November 21, 2026].
\url{https://pages.uoregon.edu/ddugger/hocolim.pdf}.
Theorem 19.4 and the following derived-colimit discussion,
p.~91; Proposition 19.9, p.~93.
Earlier shorter versions have different proposition numbers.

The Stacks Project Authors. (2026). \emph{Dold--Kan}.
Theorem 14.24.3, Tag \texttt{019G}; Section 14.24, Tag
\texttt{019D}.
\url{https://stacks.math.columbia.edu/tag/019G}.

